% ============================================================================
% AION Autonomous General Apprentice Kernel
% Compile-ready section. Requires: amsmath, amssymb, booktabs, tabularx, array, url.
% ============================================================================

\section{The Autonomous General Apprentice Kernel}
\label{sec:aion-autonomous-general-apprentice}

The AION development programme previously accumulated capability through a
sequence of internally defined phases.  Although those phases established
persistent cognition, causal learning, tool invention, outcome reflection,
private self-repair, and transfer, phase completion was not itself a mechanism
for producing increasingly general competence.  The Autonomous General
Apprentice Kernel replaces that phase-count trajectory with one persistent,
evidence-driven programme.

The internally promoted and retained procedure is shown below.
\begin{center}
\small\path{procedure_autonomous_general_apprentice_kernel_v1}
\end{center}
Its immutable owner mission is:

\begin{quote}
Become an autonomous general apprentice that can learn unfamiliar subjects,
act against real outcome authorities, transfer methods across domains, invent
missing representations and tools, retain competence, and compound
improvement.
\end{quote}

The resulting operating chain is
\begin{equation}
\begin{aligned}
\text{broad immutable mission}
&\rightarrow \text{verified-experience reconstruction}
\rightarrow \text{evidence-gap measurement} \\
&\rightarrow \text{weakest executable gap}
\rightarrow \text{specialist-contract invention} \\
&\rightarrow \text{precommitted prediction or action}
\rightarrow \text{later independent consequence} \\
&\rightarrow \text{diagnosis}
\rightarrow \text{cross-domain method transfer}
\rightarrow \text{retention} \\
&\rightarrow \text{next autonomous curriculum task}.
\end{aligned}
\label{eq:general-apprentice-loop}
\end{equation}

The kernel contains no fixed list of subjects, benchmark questions, expected
answers, or task identifiers.  Its fixed structure is limited to the
owner-approved definition of the evidence required for increasingly general
apprenticeship.  New tasks must be derived from evidence deficiencies and
available outcome authorities rather than selected from a development-authored
subject menu.

\subsection{General-Apprenticeship Evidence Model}
\label{subsec:general-apprentice-evidence-model}

The kernel tracks ten long-term evidence gates:
\begin{enumerate}
    \item open-ended learning;
    \item cross-domain transfer;
    \item learning from independently revealed consequences;
    \item autonomous curriculum construction;
    \item representation and tool invention;
    \item long-term continual learning without serious forgetting;
    \item general natural-language cognition;
    \item real planning and sustained agency;
    \item physical and social grounding; and
    \item compounding self-improvement.
\end{enumerate}
These gates specify evidence requirements, not curricula.  For gate $g$, let
$A_g$ be the set of supporting committed artifacts and $F_g$ the set of
distinct capability families represented by those artifacts.  The internal
support statistic is
\begin{equation}
S_g = 0.65\min\!\left(1,\frac{|F_g|}{4}\right)
      +0.35\min\!\left(1,\frac{|A_g|}{12}\right).
\label{eq:internal-apprentice-support}
\end{equation}
This statistic ranks internal evidence gaps only.  It is not an AGI score and
cannot mark a gate externally complete.

The next curriculum target is selected from the least-supported gate:
\begin{equation}
g^{\star}=\arg\min_{g\in\mathcal{G}}
\left(S_g,\,|A_g|,\,\operatorname{name}(g)\right).
\label{eq:weakest-apprentice-gate}
\end{equation}
The lexical term provides deterministic tie-breaking and has no effect on the
evidence score.

\subsection{Verified Experience Reconstruction and Anti-Circularity}
\label{subsec:apprentice-experience-index}

At every evidence refresh, AION reconstructs a manifest over the complete
HexCore result store.  Every artifact is SHA-256 committed and assigned one of
four epistemic roles: verified positive, verified negative, protocol-only, or
unresolved.  Positive artifacts may support retrieval and proposal formation;
negative artifacts remain available for criticism, abstention, and
calibration; protocol artifacts cannot become training truth.

The first Autonomous General Apprentice run reconstructed the evidence shown in
Table~\ref{tab:general-apprentice-corpus}.

\begin{table}[htbp]
\centering
\caption{Verified-experience reconstruction for the first apprentice run.}
\label{tab:general-apprentice-corpus}
\begin{tabular}{@{}lr@{}}
\toprule
Measure & Result \\
\midrule
SHA-256 committed artifacts indexed & 141 \\
Eligible prior verified positives & 107 \\
Preserved verified failures & 21 \\
Capability families represented & 7 \\
General-apprenticeship gates tracked & 10 \\
Cross-family method cards & 3 \\
Active procedure artifacts excluded from its evidence & 1 \\
\bottomrule
\end{tabular}
\end{table}

The active procedure is explicitly excluded from its own evidence index:
\begin{equation}
\mathcal{A}^{\mathrm{eligible}}
=\left\{a\in\mathcal{A}:a.\mathrm{procedure\_id}
\neq \mathrm{procedure\_id}_{\mathrm{active}}\right\}.
\label{eq:self-evidence-exclusion}
\end{equation}
This prevents a previously written promotion artifact from recursively
strengthening the evidence used to justify the same procedure.

\subsection{Cross-Domain Method Cards}
\label{subsec:apprentice-method-cards}

For every CAU-promoted outcome, the kernel extracts the procedure's verified
method steps, removes purely syntactic connective terms, and binds the resulting
method signature to its artifact, procedure, and capability family.  A method
$m$ becomes a transfer candidate only when
\begin{equation}
\left|\left\{f:\,m\text{ occurs in a verified artifact from family }f\right\}\right|
\geq 2.
\label{eq:method-transfer-card}
\end{equation}
This stores a method and its provenance rather than treating one domain's
answer as a universal rule.

Three cross-family method cards were recovered in the first run.  The strongest
method available to the initial curriculum task was experiment selection by
expected information gain per unit cost.  That method had independently
survived schema-revision and tool/program-induction outcomes before being
proposed for public change monitoring.

\subsection{Menu-Free Curriculum and Specialist Invention}
\label{subsec:menu-free-apprentice-curriculum}

After selecting $g^{\star}$, the kernel chooses an unused authority permitted
by the immutable mission.  Authorities are discovered from executable tools,
delayed public-outcome ledgers, and retained evidence stores.  For authority
$a$, a task contract contains
\begin{equation}
T=(I,O,A,P,X,R),
\end{equation}
where $I$ and $O$ are invented inputs and outputs, $A$ is the outcome authority,
$P$ is the precommitment, $X$ is the source-disjoint transfer requirement, and
$R$ is the restart-retention requirement.

If no permitted authority can test the selected gap, the kernel emits
\path{invent_or_acquire_outcome_authority}.  It does not substitute an internal
score.  Every generated contract requires:
\begin{itemize}
    \item a later independently verified consequence;
    \item positive transfer to a different capability family;
    \item protected backward retention;
    \item zero unsafe actions; and
    \item zero mutation of the terminal owner objective.
\end{itemize}

The first curriculum task contained no human-supplied task fields and ten
derived fields.  The kernel selected autonomous curriculum as the weakest
current evidence gate and selected the live Arena~v15 public-outcome ledger as
its authority.

\subsection{Precommitment to a Delayed Public Consequence}
\label{subsec:apprentice-delayed-public-outcome}

The first task observed five completed Arena~v15 cycles and predicted which
sources would change in the next cycle.  The prediction method used recent
change persistence with historical frequency as a fallback.  Before a new
outcome existed, AION committed
\begin{equation}
C=\operatorname{SHA256}
\left(\mathrm{task\_id},\,5,\,
\{\text{node},\text{open\_meteo\_madrid}\}\right),
\end{equation}
where
\begin{center}
\path{C = b0f0195ea5b552311aa63f667e56d71c3cd7b46cb3419349f6f40d3250f29f18}.
\end{center}

Let $c=5$ be the committed ledger cursor and $L$ the current ledger.  An outcome
is eligible only when
\begin{equation}
|L|>c,
\end{equation}
and the service binds exactly the first row $L_{c+1}$.  Rows $L_1,\ldots,L_c$
cannot be reused after the prediction is known.  For predicted change set $P$
and observed set $O$, the consequence score is the Jaccard agreement
\begin{equation}
J(P,O)=
\begin{cases}
\dfrac{|P\cap O|}{|P\cup O|}, & P\cup O\neq\varnothing,\\[6pt]
1, & P\cup O=\varnothing.
\end{cases}
\label{eq:apprentice-delayed-score}
\end{equation}
The outcome is accepted only if the public ledger also records safe execution
and a successful transaction.  Tests verify that all precommitment rows are
rejected and that the later row preserves its independent outcome hash.

\subsection{Persistent Runtime Integration}
\label{subsec:general-apprentice-runtime-integration}

The persistent service runs beside the canonical cognitive runtime and the
real-outcome learner:
\begin{center}
\small\path{backend/AION/system/aion_general_apprentice_service.py}
\end{center}
It periodically:
\begin{enumerate}
    \item rebuilds the verified-experience index;
    \item preserves the terminal mission hash;
    \item checks for a postcommitment consequence;
    \item binds and scores an eligible outcome;
    \item closes the current learning generation; and
    \item constructs the next weakest-gap curriculum task.
\end{enumerate}
The service is registered with AION Heartbeat for supervised startup and
restart.  Runtime state preserves the mission, evidence revision, method cards,
curriculum history, outcome receipts, and generation summaries.

\subsection{Promotion Results}
\label{subsec:general-apprentice-results}

\begin{table}[htbp]
\centering
\caption{Autonomous General Apprentice control-plane results.}
\label{tab:general-apprentice-results}
\begin{tabular}{@{}lr@{}}
\toprule
Measure & Result \\
\midrule
Curriculum generated without subject/task menu & Passed \\
Human-supplied fields in generated task & 0 \\
Automatically derived task fields & 10 \\
Terminal owner objective immutable & Passed \\
Missing-authority behaviour & Fail closed \\
Restart retention & Passed \\
Self-evidence exclusion & Passed \\
Unsafe actions & 0 \\
Objective mutations & 0 \\
Focused integrated regression tests & 19/19 \\
\bottomrule
\end{tabular}
\end{table}

The first execution promoted the apprentice kernel.  Subsequent identical
executions correctly returned \path{NO_CHAMPION_IMPROVEMENT} while retaining
\path{procedure_autonomous_general_apprentice_kernel_v1} as champion.  An
idempotent rerun therefore cannot create a duplicate promotion.

\subsection{Authority and Claim Boundary}
\label{subsec:general-apprentice-claim-boundary}

The promoted result establishes the persistent control plane, not completion of
the ten long-term competence gates.  The retained corpus remains predominantly
internal and development-authored.  The first live curriculum prediction is
committed but awaits a strictly later public consequence.  Hundreds or
thousands of unfamiliar tasks, multi-month retention, broad social grounding,
and sustained positive compounding have not yet been demonstrated.

Accordingly, this result supports the claim that AION now possesses a running
mechanism for evidence-driven general apprenticeship.  It does not support a
claim that AION is already an autonomous general apprentice, a generally
capable self-improving system, or AGI.  Future progress must be measured by
declining human scaffolding, positive transfer to qualitatively different
domains, lower investigation cost, preserved prior competence, and successful
outcomes under authorities AION does not control.
